\chapter[Combinatorial frontier]{The combinatorial frontier of chiral bar cohomology}
\label{app:combinatorial-frontier}

This appendix consolidates all known enumerative data arising from the chiral bar-cobar theory and organizes the computational frontier.  The sequences, generating functions, and numerical invariants below are the \emph{shadow-level} (S-level) output of the algebraic machinery developed in Parts~1--3; each entry records the theorem or conjecture from which it derives, its current verification status, and the obstruction to further computation.

\medskip
\noindent\emph{Principle.}
For a Koszul chiral algebra~$\cA$, the bar cohomology dimensions $\dim H^n(\barBgeom(\cA))$ coincide with the Hilbert function of the Koszul dual $\cA^!$ (Corollary~\ref{cor:bar-cohomology-koszul-dual}).  These dimensions are the central enumerative invariants of the theory.  Their generating functions encode, in compact form, the full structure of the bar complex: growth rates, analytic singularities, and (conjecturally) the spectral data of a finite-rank transfer operator.

\section{Master table of bar cohomology dimensions}
\label{sec:master-table}

Table~\ref{tab:frontier-bar-dims} extends Table~\ref{tab:bar-dimensions} with additional algebras and deeper commentary on the status of each entry.

\begin{table}[ht]
\centering
\caption{Bar cohomology dimensions $\dim(\cA^!)_n$ through degree~$10$.
$\bullet$\,=\,proved, $\circ$\,=\,conjectured from GF/recurrence, blank\,=\,unknown.}
\label{tab:frontier-bar-dims}
\renewcommand{\arraystretch}{1.35}
{\scriptsize
\begin{tabular}{|l|ccccccccccc|l|l|}
\hline
\textbf{Algebra $\cA$}
  & \textbf{0} & \textbf{1} & \textbf{2} & \textbf{3} & \textbf{4} & \textbf{5} & \textbf{6} & \textbf{7} & \textbf{8} & \textbf{9} & \textbf{10}
  & \textbf{Growth} & \textbf{OEIS} \\
\hline
Heisenberg ($\mathrm{rk}\,1$)
  & $\bullet 1$ & $\bullet 1$ & $\bullet 1$ & $\bullet 1$ & $\bullet 2$ & $\bullet 3$ & $\bullet 5$ & $\bullet 7$ & $\bullet 11$ & $\bullet 15$ & $\bullet 22$
  & $e^{c\sqrt{n}}$ & A000041 \\
\hline
Free fermion $\psi$
  & $\bullet 1$ & $\bullet 1$ & $\bullet 1$ & $\bullet 2$ & $\bullet 3$ & $\bullet 5$ & $\bullet 7$ & $\bullet 11$ & $\bullet 15$ & $\bullet 22$ & $\bullet 30$
  & $e^{c\sqrt{n}}$ & A000041 \\
\hline
$bc$ ghosts
  & $\bullet 1$ & $\bullet 2$ & $\bullet 3$ & $\bullet 6$ & $\bullet 13$ & $\bullet 28$ & $\bullet 59$ & $\bullet 122$ & $\bullet 249$ & $\bullet 504$ & $\bullet 1013$
  & $2^n$ & A132045 \\
\hline
$\beta\gamma$ system
  & $\bullet 1$ & $\bullet 2$ & $\bullet 4$ & $\bullet 10$ & $\bullet 26$ & $\bullet 70$ & $\bullet 192$ & $\bullet 534$ & $\bullet 1500$ & $\bullet 4246$ & $\bullet 12092$
  & $3^n / n^{1/2}$ & A025565 \\
\hline
$\widehat{\mathfrak{sl}}_2$
  & $\bullet 1$ & $\bullet 3$ & $\bullet 6$ & $\bullet 15$ & $\bullet 36$ & $\bullet 91$ & $\bullet 232$ & $\bullet 603$ & $\bullet 1585$ & $\bullet 4213$ & $\bullet 11298$
  & $3^n / n^{3/2}$ & A005043 \\
\hline
$\widehat{\mathfrak{sl}}_3$
  & $\bullet 1$ & $\bullet 8$ & $\bullet 36$ & $\bullet 204$ & $\circ 1352$ & $\circ 9892$ & $\circ 76084$ & $\circ 598592$ & $\circ 4755444$ & $\circ 37963948$ & $\circ 303961012$
  & $8^n$ & --- \\
\hline
$\mathrm{Vir}_c$
  & $\bullet 1$ & $\bullet 1$ & $\bullet 2$ & $\bullet 5$ & $\bullet 12$ & $\bullet 30$ & $\bullet 76$ & $\bullet 196$ & $\bullet 512$ & $\bullet 1353$ & $\bullet 3610$
  & $3^n / n^{3/2}$ & A002026 \\
\hline
$\mathcal{W}_3$
  & $\bullet 1$ & $\bullet 2$ & $\bullet 5$ & $\bullet 16$ & $\bullet 52$ & $\circ 171$ & $\circ 564$ & $\circ 1862$ & $\circ 6149$ & $\circ 20305$ & $\circ 67044$
  & $3.30^n$ & --- \\
\hline
$Y(\mathfrak{sl}_2)$
  & $\bullet 1$ & $\bullet 4$ & $\bullet 10$ & $\bullet 28$ & $\circ 82$ & $\circ 244$ & $\circ 730$ & $\circ 2188$ & $\circ 6562$ & $\circ 19684$ & $\circ 59050$
  & $3^n$ & --- \\
\hline
$\widehat{G}_{2,k}$
  & $\bullet 1$ & $\bullet 14$ & & & & & & & & &
  & $14^n$ (pred.) & --- \\
\hline
$\widehat{\mathfrak{sp}}_{4,k}$
  & $\bullet 1$ & $\bullet 10$ & & & & & & & & &
  & $10^n$ (pred.) & --- \\
\hline
\end{tabular}
}

\medskip
\noindent
{\footnotesize
$\bullet$\,Proved values derive from:
PBW spectral sequence collapse (Theorem~\ref{thm:spectral-sequence-collapse}),
chiral Koszulness (Theorems~\ref{thm:km-chiral-koszul},
\ref{thm:virasoro-chiral-koszul}),
and explicit computation of the bar differential rank.
$\circ$\,Conjectured values derive from fitting a rational or algebraic
generating function to the proved data points; see
Conjectures~\ref{conj:sl3-bar-gf}, \ref{conj:w3-algebraicity},
\ref{conj:yangian-bar-gf}.
For $\widehat{G}_{2,k}$ and $\widehat{\mathfrak{sp}}_{4,k}$,
only $H^0 = 1$ and $H^1 = \dim\mathfrak{g}$ are known
(the latter by Whitehead).
}
\end{table}

\begin{remark}[Reading the table]
\label{rem:reading-frontier-table}
The zeroth column ($n=0$) is always~$1$: the augmentation $\cA \to \mathbb{C}$ contributes a one-dimensional class in bar degree zero.  The first column ($n=1$) equals $\dim\overline{\cA}_{[1]}$, the dimension of the weight-$1$ augmentation ideal.  For Kac--Moody algebras $\widehat{\mathfrak{g}}_k$, this is $\dim\mathfrak{g}$; for the Virasoro algebra, it is~$1$ (the single generator~$T$); for $\mathcal{W}_3$, it is~$2$ (generators $T$ and~$W$).

\emph{The data is level-independent.}  For every algebra in the table, the bar cohomology dimensions are independent of the level~$k$ (or central charge~$c$) at generic values.  This is a deep consequence of the PBW spectral sequence structure: the $E_1$ page is a Lie algebra cohomology computation that depends only on the underlying finite-dimensional data $(\mathfrak{g}, \dim\mathfrak{g}, h^\vee)$, not on the level.  The level enters only through the curvature term $\omega_g$ at genus $g \geq 1$, which affects the genus expansion (Theorem~\ref{thm:genus-universality}) but not the bar cohomology dimensions themselves.
\end{remark}

\section{Named combinatorial sequences}
\label{sec:named-sequences}

Six classical combinatorial sequences appear as bar cohomology dimensions.  For each, we record the sequence definition, the mechanism by which it arises, and the OEIS reference.

\subsection{Partition numbers (Heisenberg, free fermion)}
\label{subsec:partition-numbers}

The partition function $p(n) = |\{\lambda \vdash n\}|$ (OEIS~A000041) enumerates integer partitions.

\begin{center}
\renewcommand{\arraystretch}{1.2}
\begin{tabular}{c|cccccccccc}
$n$ & 0 & 1 & 2 & 3 & 4 & 5 & 6 & 7 & 8 & 9 \\
\hline
$p(n)$ & 1 & 1 & 2 & 3 & 5 & 7 & 11 & 15 & 22 & 30
\end{tabular}
\end{center}

\emph{How it appears.}
The Heisenberg algebra $\mathcal{H}_k$ has $\dim(\mathcal{H}_k^!)_n = p(n-2)$ for $n \geq 2$ and $\dim(\mathcal{H}_k^!)_0 = \dim(\mathcal{H}_k^!)_1 = 1$ (Theorem~\ref{thm:heisenberg-bar}).  The free fermion $\psi$ has $\dim(\psi^!)_n = p(n-1)$.  Both identifications follow from the structure of the augmentation ideal: the Heisenberg is generated by a single bosonic field $\alpha(z)$ of weight~$1$, whose descendants $\alpha_{-n}$ ($n \geq 1$) form an abelian Lie algebra isomorphic to $\mathbb{C}[t^{-1}]t^{-1}$.  The bar complex reduces to the exterior algebra of this abelian Lie algebra, graded by conformal weight, and the weight-$h$ component has dimension~$p(h)$.

\emph{Generating function.}
The partition function has the celebrated Euler product:
\[
\sum_{n \geq 0} p(n)\, q^n = \prod_{k=1}^{\infty} \frac{1}{1 - q^k}
= \frac{1}{\eta(q) \cdot q^{-1/24}},
\]
where $\eta(q) = q^{1/24}\prod_{n \geq 1}(1-q^n)$ is the Dedekind eta function.  This is the genus-$1$ partition function of the free boson (Theorem~\ref{thm:genus1-heisenberg}).  The transcendence of this generating function (it is not algebraic over $\mathbb{Q}(q)$) reflects the abelian, non-interacting nature of the Heisenberg algebra.

\emph{Asymptotics.}
Hardy--Ramanujan: $p(n) \sim e^{\pi\sqrt{2n/3}}/(4\sqrt{3}\,n)$, sub-exponential growth.  This characterizes free fields among Koszul chiral algebras (Corollary~\ref{cor:subexp-free-field}).

\subsection{Riordan numbers ($\widehat{\mathfrak{sl}}_2$)}
\label{subsec:riordan}

The Riordan numbers $R(n)$ (OEIS~A005043) satisfy the recurrence
\begin{equation}\label{eq:riordan-recurrence}
(n+1)\,R(n) = (n-1)\bigl(2R(n-1) + 3R(n-2)\bigr),
\qquad R(0) = 1,\; R(1) = 0,
\end{equation}
with generating function
\begin{equation}\label{eq:riordan-gf}
\sum_{n \geq 0} R(n)\,x^n
= \frac{1 + x - \sqrt{1-2x-3x^2}}{2x(1+x)}\,.
\end{equation}

\begin{center}
\renewcommand{\arraystretch}{1.2}
\begin{tabular}{c|cccccccccccc}
$n$ & 0 & 1 & 2 & 3 & 4 & 5 & 6 & 7 & 8 & 9 & 10 & 11 \\
\hline
$R(n)$ & 1 & 0 & 1 & 1 & 3 & 6 & 15 & 36 & 91 & 232 & 603 & 1585
\end{tabular}
\end{center}

\emph{How it appears.}
$\dim(\widehat{\mathfrak{sl}}_{2,k}^!)_n = R(n+3)$ for all $n \geq 0$ (Theorem~\ref{thm:universal-kac-moody-koszul}, Corollary~\ref{cor:bar-cohomology-koszul-dual}).  This is the first instance of a \emph{named} combinatorial sequence emerging from chiral bar cohomology.  The shift by~$3$ arises because the augmentation ideal of $\widehat{\mathfrak{sl}}_2$ starts at conformal weight~$1$, and the first three Riordan numbers $R(0) = 1$, $R(1) = 0$, $R(2) = 1$ correspond to the trivial bar degrees.

\emph{Combinatorial interpretation.}
Riordan numbers count several classical objects~\cite{BernhartRiordan}: Motzkin paths of length~$n$ that never touch the $x$-axis except at endpoints; non-crossing partitions of $\{1, \ldots, n\}$ without singleton blocks; and certain classes of restricted lattice paths.  The appearance of Riordan numbers in chiral bar cohomology suggests that the bar complex of $\widehat{\mathfrak{sl}}_2$ admits a combinatorial model in terms of non-crossing structures --- an analogue of the Motzkin path model for the Virasoro algebra (Theorem~\ref{conj:motzkin-path-model}).  No such model has been constructed.

\emph{Asymptotics.}
$R(n) \sim (3/4)\sqrt{3/\pi}\, n^{-3/2}\, 3^n$ (from the square-root singularity at $x = 1/3$).  The growth rate~$3^n = (\dim\mathfrak{sl}_2)^n$ is sharp: the chain-space bound $\dim\bar{B}^n \leq (\dim\mathfrak{g})^n \cdot (n-1)!$ shows that the exponential growth rate of bar cohomology cannot exceed $\dim\mathfrak{g}$ (Corollary~\ref{cor:growth-rate-dimg}), and for $\widehat{\mathfrak{sl}}_2$ this bound is saturated.

\subsection{Motzkin numbers and their differences (Virasoro)}
\label{subsec:motzkin}

The Motzkin numbers $M(n)$ (OEIS~A001006) count lattice paths from $(0,0)$ to $(n,0)$ using up steps $(1,1)$, down steps $(1,-1)$, and flat steps $(1,0)$, staying at height $\geq 0$.

\begin{center}
\renewcommand{\arraystretch}{1.2}
\begin{tabular}{c|cccccccccccc}
$n$ & 0 & 1 & 2 & 3 & 4 & 5 & 6 & 7 & 8 & 9 & 10 & 11 \\
\hline
$M(n)$ & 1 & 1 & 2 & 4 & 9 & 21 & 51 & 127 & 323 & 835 & 2188 & 5798 \\
$M(n{+}1) - M(n)$ & 0 & 1 & 2 & 5 & 12 & 30 & 76 & 196 & 512 & 1353 & 3610 & 9694
\end{tabular}
\end{center}

\emph{How it appears.}
$\dim(\mathrm{Vir}_c^!)_n = M(n+1) - M(n)$ for all $n \geq 1$ (Theorem~\ref{thm:spectral-sequence-collapse}).  The differences $M(n+1) - M(n)$ (OEIS~A002026) count $U$-initial Motzkin paths of length~$n+1$, i.e., Motzkin paths whose first step is an up step.  This count has a direct representation-theoretic meaning: the bijection of Theorem~\ref{conj:motzkin-path-model} associates to each basis element of $H^n(\barBgeom(\mathrm{Vir}_c))$ a $U$-initial Motzkin path, with the three step types corresponding to the three singular terms in the Virasoro OPE $T(z)T(w)$.

\emph{Generating function.}
\[
\sum_{n \geq 0} M(n)\,x^n
= \frac{1 - x - \sqrt{1-2x-3x^2}}{2x^2}\,,
\]
sharing the discriminant $\Delta(x) = 1-2x-3x^2 = (1-3x)(1+x)$ with the Riordan generating function.  This shared discriminant is the signature of the $\mathfrak{sl}_2$ DS family (Theorem~\ref{thm:ds-bar-gf-discriminant}).

\emph{The Motzkin recurrence.}
$(n+2)\,M(n) = (2n+1)\,M(n-1) + 3(n-1)\,M(n-2)$, equivalent to the algebraic equation $x^2 M^2 + (x-1)M + 1 = 0$.  The $P$-recursive structure reflects the Koszul property: bar cohomology dimensions of Koszul chiral algebras satisfy holonomic recurrences.

\subsection{Catalan numbers (A-infinity, associahedra)}
\label{subsec:catalan}

The Catalan numbers $C_n = \frac{1}{n+1}\binom{2n}{n}$ (OEIS~A000108) appear not as bar cohomology dimensions directly, but as the combinatorial skeleton underlying the homotopy transfer machinery.

\begin{center}
\renewcommand{\arraystretch}{1.2}
\begin{tabular}{c|cccccccccc}
$n$ & 0 & 1 & 2 & 3 & 4 & 5 & 6 & 7 & 8 & 9 \\
\hline
$C_n$ & 1 & 1 & 2 & 5 & 14 & 42 & 132 & 429 & 1430 & 4862
\end{tabular}
\end{center}

\emph{How they appear.}

\begin{enumerate}[label=(\roman*)]
\item \emph{Planar rooted trees.}  The number of planar rooted trees with~$n$ internal nodes is $C_{n-1}$.  These trees index the terms of the $A_\infty$ structure maps $m_n$ in the homotopy transfer theorem (Appendix~\ref{app:homotopy-transfer}, Theorem~\ref{thm:htt}).  When one transfers the bar complex structure to a minimal model via a strong deformation retract, each $m_n$ is a sum over $C_{n-1}$ planar trees.

\item \emph{Associahedra.}  The associahedron $K_n$ has $C_{n-1}$ vertices (maximal bracketings of an $(n+1)$-fold product).  The FM compactification $\overline{C}_n(\mathbb{P}^1) \cong K_n$ of the genus-$0$ configuration space is isomorphic to the associahedron as a stratified space (Theorem~\ref{thm:fm-associahedron}).  The Betti numbers of $K_n$ are related to Narayana numbers, and the $f$-vector encodes the bar spectral sequence page dimensions.

\item \emph{Virasoro at low degree.}  The first few Virasoro bar cohomology dimensions $(1, 2, 5, 12, 30)$ begin with a Catalan-adjacent pattern: $C_1 = 1$, $C_2 = 2$, $C_3 = 5$.  The agreement is not coincidental but breaks at $n = 4$ ($C_4 = 14 \neq 12$), reflecting the distinction between the full Catalan enumeration and the Motzkin-difference count.
\end{enumerate}

\subsection{Bernoulli numbers (genus expansion)}
\label{subsec:bernoulli}

The Bernoulli numbers $B_{2g}$ (OEIS~A027642 for denominators) govern the genus expansion of every Koszul chiral algebra.

\begin{center}
\renewcommand{\arraystretch}{1.2}
\begin{tabular}{c|cccccccccc}
$g$ & 1 & 2 & 3 & 4 & 5 & 6 & 7 & 8 & 9 & 10 \\
\hline
$|B_{2g}|$ & $\tfrac{1}{6}$ & $\tfrac{1}{30}$ & $\tfrac{1}{42}$ & $\tfrac{1}{30}$ & $\tfrac{5}{66}$ & $\tfrac{691}{2730}$ & $\tfrac{7}{6}$ & $\tfrac{3617}{510}$ & $\tfrac{43867}{798}$ & $\tfrac{174611}{330}$
\end{tabular}
\end{center}

\emph{How they appear.}
The Faber--Pandharipande coefficients
\[
\lambda_g^{FP} = \frac{2^{2g-1}-1}{2^{2g-1}} \cdot \frac{|B_{2g}|}{(2g)!}
\]
are the universal genus-$g$ obstruction coefficients (Proposition~\ref{prop:fp-coefficients}).  For any Koszul chiral algebra~$\cA$, the genus-$g$ free energy factors as $F_g(\cA) = \kappa(\cA) \cdot \lambda_g^{FP}$ (Theorem~\ref{thm:genus-universality}).  The generating function
\[
\sum_{g \geq 1} \lambda_g^{FP}\, x^{2g}
= \frac{x/2}{\sin(x/2)} - 1
= \hat{A}(x) - 1,
\]
where $\hat{A}(x) = x/2 / \sinh(x/2)$ is the $\hat{A}$-genus, is the even part of the Todd class.  The appearance of Bernoulli numbers via the $\hat{A}$-genus is a manifestation of the index-theoretic nature of genus corrections: the bar complex curvature at genus~$g$ is controlled by the Hodge bundle on $\overline{\mathcal{M}}_g$, and the Faber--Pandharipande formula is an intersection-theoretic avatar of the Hirzebruch--Riemann--Roch theorem.

\emph{Asymptotics.}
$\lambda_g^{FP} \sim 2/(2\pi)^{2g}$, super-exponential decay.  This guarantees convergence of the genus expansion $\sum_g F_g(\cA)\,x^{2g}$ for all finite~$\kappa$ (Proposition~\ref{prop:bivariate-gf}).

\subsection{Stirling numbers (configuration spaces)}
\label{subsec:stirling}

The unsigned Stirling numbers of the first kind $|s(n,k)| = e_k(1, 2, \ldots, n-1)$ (elementary symmetric polynomials evaluated at $1, \ldots, n-1$) appear as Betti numbers of configuration spaces.

\emph{How they appear.}
The Poincar\'e polynomial of the configuration space $C_n(\mathbb{C})$ is
\[
P_t(C_n(\mathbb{C})) = \prod_{j=1}^{n-1}(1+jt)
= \sum_{k=0}^{n-1} e_k(1, 2, \ldots, n-1)\, t^k,
\]
so $\dim H^k(C_n(\mathbb{C})) = e_k(1, \ldots, n-1)$.  In particular:
\begin{itemize}
\item $\dim H^0 = 1$ (connected);
\item $\dim H^1 = \binom{n}{2}$ (Arnold generators $\eta_{ij}$);
\item $\dim H^{n-1} = (n-1)!$ (top degree; equals $\dim\mathrm{OS}^{n-1}(n)$).
\end{itemize}
The top-degree dimension $(n-1)!$ directly enters the chain-group formula: $\dim\bar{B}^n(\widehat{\mathfrak{g}}_k) = (\dim\mathfrak{g})^n \cdot (n-1)!$ (Remark~\ref{rem:bar-dims-level-independent}).  The factorial growth of the chain space, contrasted with the exponential growth of bar cohomology, measures the efficiency of the bar differential: cancellation accounts for $(n-1)!/(\dim\mathfrak{g})^n$ of the chain space at leading order.


\section{Generating functions: the algebraic hierarchy}
\label{sec:gf-hierarchy}

The generating functions of bar cohomology dimensions fall into a strict hierarchy of algebraic complexity.  This hierarchy reflects the interaction structure of the underlying chiral algebra.

\subsection{Transcendental (free fields)}
\label{subsec:gf-transcendental}

\begin{align}
P_{\mathcal{H}}(q) &= \prod_{k \geq 1} \frac{1}{1-q^k}
&&\text{(Heisenberg; weight-graded)} \\
P_{\psi}(q) &= \prod_{k \geq 1} (1+q^k)
&&\text{(free fermion; weight-graded)}
\end{align}

These are modular objects (related to $1/\eta(\tau)$ and $\eta(2\tau)/\eta(\tau)$ respectively).  Their transcendence over $\mathbb{Q}(q)$ is a consequence of the Nesterenko--Bertrand theorem on modular transcendence.  No polynomial differential equation over $\mathbb{Q}(q)$ is satisfied.  This reflects the absence of nontrivial OPE poles: the bar differential is zero, and the bar cohomology is the full chain space.

\emph{Conformal-weight-collapsed series.}  If one forgets the conformal weight grading and counts only by bar degree, the Heisenberg dimensions become $\dim(\mathcal{H}^!)_n \in \{1, 1, 1, 2, 3, 5, 7, 11, \ldots\} = p(n-2)$ (Table~\ref{tab:frontier-bar-dims}).

\subsection{Rational (bc ghosts, Yangian)}
\label{subsec:gf-rational}

\begin{align}
P_{bc}(x) &= \sum_{n \geq 1}(2^n - n + 1)\,x^n
&&\text{(closed form: $\frac{2x}{1-2x} - \frac{x}{(1-x)^2} + \frac{x}{1-x}$)}
\label{eq:bc-gf}
\end{align}

The $bc$ ghost system has a rational generating function because its OPE is purely fermionic with a single pole of order~$1$: the bar differential is maximally sparse, and the cohomology computation reduces to a subtraction of two geometric series.

\emph{Yangian (conjectured).}
$P_{Y(\mathfrak{sl}_2)}(x) = (1-3x^2)/((1-x)(1-3x))$ (Conjecture~\ref{conj:yangian-bar-gf}).  If confirmed, this would be the only \emph{interacting} algebra with a rational bar generating function.  The partial fraction decomposition $-1 + 1/(1-x) + 1/(1-3x)$ suggests a K\"unneth-type splitting into a ``$\mathfrak{gl}_1$ center'' contribution ($1/(1-x)$, polynomial growth) and an ``$\mathfrak{sl}_2$ factor'' contribution ($1/(1-3x)$, exponential growth $3^n$).  This splitting is consistent with the RTT presentation of $Y(\mathfrak{sl}_2)$ (Theorem~\ref{thm:yangian-bar-rtt}).

\emph{Why rationality is surprising.}
For $E_\infty$-chiral algebras (Kac--Moody, Virasoro), the bar generating function is algebraic of degree~$2$ over $\mathbb{Q}(x)$ (Corollary~\ref{cor:algebraicity-koszul}).  The Yangian is $E_1$-chiral, not $E_\infty$, and the rationality of its bar GF may reflect the stronger algebraic constraints of the $E_1$ structure (full associativity of the chiral product, not just homotopy associativity).

\subsection{Algebraic of degree 2 ($\mathfrak{sl}_2$ family)}
\label{subsec:gf-algebraic}

The three algebras in the $\mathfrak{sl}_2$ DS orbit share a common discriminant:

\begin{center}
\renewcommand{\arraystretch}{1.4}
\begin{tabular}{lll}
\toprule
\textbf{Algebra} & \textbf{Generating function} & \textbf{Singularity type} \\
\midrule
$\widehat{\mathfrak{sl}}_2$ & $(1+x-\sqrt{\Delta})\,/\,(2x(1+x))$ & $\Delta^{1/2}$ \\
$\mathrm{Vir}_c$ & $4x\,/\,(1-x+\sqrt{\Delta})^2$ & $\Delta^{1/2}$ \\
$\beta\gamma$ & $(1+x)\,/\,\sqrt{\Delta}$ & $\Delta^{-1/2}$ \\
\bottomrule
\end{tabular}
\end{center}
where $\Delta(x) = (1-3x)(1+x) = 1-2x-3x^2$.

\emph{Why the same discriminant.}
Drinfeld--Sokolov reduction $\widehat{\mathfrak{sl}}_2 \twoheadrightarrow \mathrm{Vir}_c$ preserves the discriminant (Theorem~\ref{thm:ds-bar-gf-discriminant}).  The Wakimoto representation $\widehat{\mathfrak{sl}}_2 \hookrightarrow \beta\gamma \otimes \mathcal{H}$ transmits the discriminant to the $\beta\gamma$ factor (the Heisenberg tensor factor contributes only sub-exponential corrections).  The invariance of the discriminant under DS reduction is the first proved instance of a general principle (Proposition~\ref{prop:discriminant-characteristic}).

\emph{The branch points.}
$x = 1/3$: the dominant singularity, giving exponential growth $3^n = (\dim\mathfrak{sl}_2)^n$.  $x = -1$: the subdominant singularity, giving alternating corrections $(-1)^n$.  The two eigenvalues of the Picard--Fuchs companion matrix $T_{\mathrm{rec}} = \bigl(\begin{smallmatrix} 2 & 3 \\ 1 & 0 \end{smallmatrix}\bigr)$ are precisely $3$ and $-1$ (Proposition~\ref{prop:hred-sl2}).

\subsection{Conjectured rational ($\mathfrak{sl}_3$ family)}
\label{subsec:gf-sl3-family}

The $\mathfrak{sl}_3$ DS orbit is conjectured to have \emph{rational} (not algebraic) generating functions sharing the denominator factor $(1-3x-x^2)$:

\begin{center}
\renewcommand{\arraystretch}{1.4}
\begin{tabular}{lll}
\toprule
\textbf{Algebra} & \textbf{Conjectured GF} & \textbf{Denominator} \\
\midrule
$\widehat{\mathfrak{sl}}_3$ & $4x(2-13x-2x^2)\,/\,\bigl((1-8x)(1-3x-x^2)\bigr)$ & $(1-8x)(1-3x-x^2)$ \\
$\mathcal{W}_3$ & $x(2-3x)\,/\,\bigl((1-x)(1-3x-x^2)\bigr)$ & $(1-x)(1-3x-x^2)$ \\
\bottomrule
\end{tabular}
\end{center}

\emph{Characteristic polynomial.}
The recurrence $a(n) = 11a(n{-}1) - 23a(n{-}2) - 8a(n{-}3)$ for $\widehat{\mathfrak{sl}}_3$ has characteristic polynomial $t^3 - 11t^2 + 23t + 8 = (t-8)(t^2-3t-1)$.  The factorization mirrors the $\mathfrak{sl}_2$ case: the dominant root $t = 8 = \dim\mathfrak{sl}_3$ corresponds to the growth pole, while the quadratic factor $t^2 - 3t - 1$ is the reciprocal polynomial of the DS-invariant discriminant $1 - 3x - x^2$.  Its roots are $(3 \pm \sqrt{13})/2$.

\emph{Why rationality, not algebraicity?}
The shift from algebraic (degree~$2$) to rational is unexpected.  A speculative explanation: the $\mathfrak{sl}_2$ Koszul dual Hilbert series satisfies a quadratic algebraic equation because $\operatorname{rank}(\mathfrak{sl}_2) = 1$, and the single Casimir operator gives a degree-$2$ recurrence.  For $\mathfrak{sl}_3$, $\operatorname{rank} = 2$, and the two independent Casimir operators give a degree-$3$ recurrence.  If the algebraic equation ``linearizes'' --- i.e., if the bar cohomology dimensions satisfy a \emph{linear} recurrence with constant coefficients, rather than one with polynomial coefficients --- then the generating function is rational.  This linearization is not understood.

\emph{The rank-plus-one pattern.}
For $\mathfrak{sl}_2$: recurrence depth $= 2 = \operatorname{rank} + 1$.
For $\mathfrak{sl}_3$ (conjectured): recurrence depth $= 3 = \operatorname{rank} + 1$.
If this pattern persists, $\widehat{\mathfrak{sl}}_4$ should satisfy a depth-$4$ recurrence, and the DS-invariant discriminant factor should be a cubic polynomial in~$x$.  See Conjecture~\ref{conj:discriminant-ks-operator} and Remark~\ref{rem:rank-plus-one}.


\section{Discriminant families and DS orbits}
\label{sec:discriminant-families}

The discriminant of the bar generating function is a Drinfeld--Sokolov invariant (Theorem~\ref{thm:ds-bar-gf-discriminant}).  This organizes the computed examples into families.

\begin{center}
\renewcommand{\arraystretch}{1.4}
\begin{tabular}{llll}
\toprule
\textbf{Family} & \textbf{DS-invariant factor} & \textbf{Roots} & \textbf{Members} \\
\midrule
$\mathfrak{sl}_2$ & $(1-3x)(1+x)$ & $1/3,\; -1$ & $\widehat{\mathfrak{sl}}_2$, $\mathrm{Vir}$, $\beta\gamma$ \\
$\mathfrak{sl}_3$ (conj.) & $1-3x-x^2$ & $\frac{-3\pm\sqrt{13}}{2}$ & $\widehat{\mathfrak{sl}}_3$, $\mathcal{W}_3$ \\
abelian & $1$ (no branch) & --- & $\mathcal{H}$, $\psi$ \\
\bottomrule
\end{tabular}
\end{center}

\begin{remark}[The number $\sqrt{13}$]
The irrational number $\sqrt{13}$ enters the theory through the $\mathfrak{sl}_3$ discriminant $1-3x-x^2 = 0$, giving roots at $x = (-3 \pm \sqrt{13})/2$.  This is \emph{not} a coincidence of small numbers: $13 = 3^2 + 4 = 3^2 + 4 \cdot 1^2$, and the discriminant of the quadratic $x^2 - 3x - 1$ is $9 + 4 = 13$.  Whether~$\sqrt{13}$ has deeper representation-theoretic meaning (e.g., as an eigenvalue of a natural operator on the $\mathfrak{sl}_3$ Cartan subalgebra) is unknown.  For comparison, the $\mathfrak{sl}_2$ family has $\sqrt{1+8} = 3$ (a perfect square), so the $\mathfrak{sl}_2$ discriminant splits over~$\mathbb{Q}$.
\end{remark}

\begin{remark}[Predicted families]
Conjecture~\ref{conj:non-simply-laced-discriminant} predicts additional DS families:
\begin{itemize}
\item $\mathfrak{sl}_4$ family: $\widehat{\mathfrak{sl}}_4$ and $\mathcal{W}_4$, sharing a degree-$3$ DS-invariant discriminant (rank~$3$ plus one growth pole).  Growth rate $15^n = (\dim\mathfrak{sl}_4)^n$.
\item $G_2$ family: $\widehat{G}_2$ and $\mathcal{W}(G_2)$.  Growth rate $14^n$.  The DS-invariant discriminant is predicted to have degree $\operatorname{rank}(G_2) + 1 = 3$.
\item $\mathfrak{sp}_4$ family: non-simply-laced.  $h^\vee = 3 \neq h = 4$.  The dual Coxeter discrepancy should manifest in the discriminant structure.
\end{itemize}
None of these have any bar cohomology data beyond $H^1 = \dim\mathfrak{g}$.
\end{remark}


\section{Growth rates and singularity analysis}
\label{sec:growth-rates}

The exponential growth rate of $\dim(\cA^!)_n$ is determined by the dominant singularity of the generating function.  By analytic combinatorics (Flajolet--Sedgewick~\cite{FlajoletSedgewick}), the growth is $\alpha^n \cdot n^{\beta}$ where $\alpha = 1/R$ ($R$ = radius of convergence) and the exponent~$\beta$ depends on the singularity type.

\begin{center}
\renewcommand{\arraystretch}{1.4}
\begin{tabular}{lccll}
\toprule
\textbf{Algebra} & $\alpha$ & $\beta$ & \textbf{Singularity type} & \textbf{Status} \\
\midrule
$\mathcal{H}$, $\psi$ & --- & --- & essential (modular) & proved \\
$bc$ & $2$ & $0$ & rational pole & proved \\
$\beta\gamma$ & $3$ & $-1/2$ & $\Delta^{-1/2}$ (algebraic) & proved \\
$\widehat{\mathfrak{sl}}_2$ & $3$ & $-3/2$ & $\Delta^{1/2}$ (algebraic) & proved \\
$\mathrm{Vir}_c$ & $3$ & $-3/2$ & $\Delta^{1/2}$ (algebraic) & proved \\
$\widehat{\mathfrak{sl}}_3$ & $8$ & $0$ & rational pole & conjectured \\
$\mathcal{W}_3$ & $\frac{3+\sqrt{13}}{2}$ & $0$ & rational pole & conjectured \\
$Y(\mathfrak{sl}_2)$ & $3$ & $0$ & rational pole & conjectured \\
$\widehat{G}_2$ & $14$ & ? & predicted & unknown \\
\bottomrule
\end{tabular}
\end{center}

\emph{The dichotomy $\beta = -3/2$ vs.\ $\beta = 0$.}
The proved cases ($\widehat{\mathfrak{sl}}_2$, $\mathrm{Vir}$) have sub-exponential correction $n^{-3/2}$, characteristic of square-root algebraic singularities.  The conjectured cases ($\widehat{\mathfrak{sl}}_3$, $\mathcal{W}_3$, $Y(\mathfrak{sl}_2)$) have $\beta = 0$, characteristic of simple rational poles.  If confirmed, this would mean that the shift from rank~$1$ ($\mathfrak{sl}_2$) to rank~$2$ ($\mathfrak{sl}_3$) changes the analytic nature of the generating function from algebraic to rational --- a qualitative transition in combinatorial complexity.

\emph{Universality of the growth rate $3$.}
The number~$3 = \dim\mathfrak{sl}_2$ appears as the growth rate for $\widehat{\mathfrak{sl}}_2$, $\mathrm{Vir}$, $\beta\gamma$, \emph{and} conjecturally $Y(\mathfrak{sl}_2)$.  The first three share a DS orbit; the Yangian does not.  The coincidence $\alpha_{Y(\mathfrak{sl}_2)} = 3 = \dim\mathfrak{sl}_2$ for the Yangian is explained by the RTT presentation: the generating matrices are $2 \times 2$, giving a $4$-dimensional space of generators ($= 2^2$), but the growth is controlled by the $\mathfrak{sl}_2$ factor after modding out the center.


\section{The modular characteristic and genus expansion}
\label{sec:modular-characteristic}

The modular characteristic $\kappa(\cA)$ is a scalar invariant that controls all genus corrections (Theorem~\ref{thm:modular-characteristic}).  Its computation for each algebra gives the following table.

\begin{center}
\renewcommand{\arraystretch}{1.4}
\begin{tabular}{lll}
\toprule
\textbf{Algebra} & $\kappa(\cA)$ & \textbf{Complementary partner} \\
\midrule
Heisenberg $\mathcal{H}_k$ & $k$ & $\mathcal{H}_{-k}$ \quad ($\kappa + \kappa' = 0$) \\
$\widehat{\mathfrak{sl}}_{2,k}$ & $\frac{3(k+2)}{4}$ & $\widehat{\mathfrak{sl}}_{2,-k-4}$ \quad ($\kappa + \kappa' = 0$) \\
$\widehat{\mathfrak{sl}}_{3,k}$ & $\frac{4(k+3)}{3}$ & $\widehat{\mathfrak{sl}}_{3,-k-6}$ \quad ($\kappa + \kappa' = 0$) \\
$\mathrm{Vir}_c$ & $c/2$ & $\mathrm{Vir}_{26-c}$ \quad ($\kappa + \kappa' = 13$) \\
$\mathcal{W}_3$ (central charge $c$) & $c/2$ & $\mathcal{W}_3^{100-c}$ \quad ($\kappa + \kappa' = 50$) \\
General $\widehat{\mathfrak{g}}_k$ & $\frac{(k+h^\vee)\dim\mathfrak{g}}{2h^\vee}$ & $\widehat{\mathfrak{g}}_{-k-2h^\vee}$ \quad ($\kappa + \kappa' = 0$) \\
\bottomrule
\end{tabular}
\end{center}

\emph{The genus expansion formula.}
For any algebra with modular characteristic $\kappa$, the bivariate generating function in the genus variable~$x$ is:
\[
\sum_{g \geq 1} F_g(\cA)\, x^{2g}
= \kappa(\cA) \cdot \left(\frac{x/2}{\sin(x/2)} - 1\right).
\]
This is the product of an algebra-dependent scalar ($\kappa$) and a universal transcendental function ($\hat{A}(x) - 1$).  The universal function is generated by Bernoulli numbers; the algebra-dependent factor is a rational function of the level.

\emph{Critical level.}
At $k = -h^\vee$ (the critical level), $\kappa = 0$ and every genus correction vanishes.  This is the level at which the Feigin--Frenkel center $\mathfrak{z}(\widehat{\mathfrak{g}})$ appears, and the geometric Langlands correspondence operates.

\emph{Complementarity arithmetic.}
For Kac--Moody algebras, complementarity gives $\kappa + \kappa' = 0$ (exact cancellation under $k \mapsto -k - 2h^\vee$).  For W-algebras obtained by DS reduction, the complementarity sum is nonzero and equals a characteristic constant of the DS orbit: $13$ for $\mathrm{Vir}$ (= half the bosonic string central charge), $50$ for $\mathcal{W}_3$ (= half the $W_3$ string central charge $100$).


\section{q-Series and modular structures}
\label{sec:q-series}

Several generating functions in the theory are (quasi-)modular forms or have modular transformation properties.

\subsection{Dedekind eta and partition functions}

The genus-$1$ partition function of the Heisenberg algebra is $Z_1(\mathcal{H}_k | \tau) = 1/\eta(\tau)$, where $\eta(\tau) = q^{1/24}\prod_{n \geq 1}(1-q^n)$ and $q = e^{2\pi i \tau}$.  The coefficients of $1/\eta$ are partition numbers.  For Kac--Moody algebras,
\[
Z_1(\widehat{\mathfrak{sl}}_{2,k} | \tau)
= \frac{\chi(\mathfrak{sl}_2)}{\eta(\tau)^3}
\cdot \bigl(1 + (k+2) E_2(\tau) + O(E_4)\bigr),
\]
where $\eta^3$ reflects $\dim\mathfrak{sl}_2 = 3$ and the correction involves the quasi-modular Eisenstein series~$E_2$.

\subsection{Eisenstein series and the holomorphic anomaly}

The weight-$2$ quasi-modular Eisenstein series $E_2(\tau) = 1 - 24\sum_{n \geq 1}\sigma_1(n)q^n$ is \emph{not} a modular form: it transforms with an inhomogeneous correction term.  This anomalous transformation is precisely the genus-$1$ holomorphic anomaly of the chiral bar complex.  The genus-$1$ curvature of $\widehat{\mathfrak{sl}}_{2,k}$ (Theorem~\ref{thm:sl2-genus1-curvature}) is a B-cycle monodromy integral whose $\tau$-dependence is controlled by $E_2(\tau)$.

\subsection{Minimal model characters}

The characters of the Virasoro minimal models $M(p,q)$ are:
\[
\chi_{r,s}(\tau)
= \frac{q^{h_{r,s} - c/24}}{\eta(\tau)}
\sum_{n \in \mathbb{Z}}
\bigl(q^{pqn^2 + (pr-qs)n} - q^{pqn^2 + (pr+qs)n}\bigr),
\]
with $c = 1 - 6(p-q)^2/(pq)$ and $h_{r,s} = ((pr-qs)^2 - (p-q)^2)/(4pq)$.  For the Ising model ($c = 1/2$, i.e., $M(3,4)$), the characters take the elegant form:
\[
\chi_{\mathbb{I}} = \tfrac{1}{2}\bigl(\sqrt{\theta_3/\eta} + \sqrt{\theta_4/\eta}\bigr),
\qquad
\chi_\sigma = \tfrac{1}{\sqrt{2}}\sqrt{\theta_2/\eta},
\qquad
\chi_\varepsilon = \tfrac{1}{2}\bigl(\sqrt{\theta_3/\eta} - \sqrt{\theta_4/\eta}\bigr),
\]
in terms of Jacobi theta functions.

\subsection{Lattice theta series}

For a positive-definite even lattice~$\Lambda$, the theta series $\Theta_\Lambda(\tau) = \sum_{\alpha \in \Lambda} q^{\langle\alpha,\alpha\rangle/2}$ is a modular form of weight $\mathrm{rank}(\Lambda)/2$ for a congruence subgroup.  The lattice vertex algebra $V_\Lambda$ has partition function $Z_1(V_\Lambda | \tau) = \Theta_\Lambda(\tau) / \eta(\tau)^{\mathrm{rank}(\Lambda)}$.


\section{Open problems and the computational frontier}
\label{sec:open-problems}

We organize the remaining computational challenges by estimated difficulty.

\subsection{Near-term targets (single new bar cohomology value)}

\begin{enumerate}[label=\textbf{NF-\arabic*.}, leftmargin=3em]
\item\label{nf:sl3-h4}
\emph{$\dim H^4(\barBgeom(\widehat{\mathfrak{sl}}_{3,k})) \stackrel{?}{=} 1352$.}
This is the single most informative computation on the frontier.  The value $1352$ would confirm the rational GF (Conjecture~\ref{conj:sl3-bar-gf}) and the shared discriminant with $\mathcal{W}_3$; any other value would refute the rational form and require reconsidering the algebraicity class.

\emph{Obstruction}: the computation requires the rank of a $24576 \times 786432$ matrix over $\mathbb{Q}$ (or a suitable finite field $\mathbb{F}_p$).  This is a $\sim 20$\,GB dense matrix, beyond naive methods but within reach of sparse or modular arithmetic.

\item\label{nf:w3-h5}
\emph{$\dim H^5(\barBgeom(\mathcal{W}_3^k)) \stackrel{?}{=} 171$.}
Confirms the $\mathcal{W}_3$ recurrence (Conjecture~\ref{conj:w3-algebraicity}).  The computation is smaller than~\ref{nf:sl3-h4} because $\mathcal{W}_3$ has only two generators (compared to $8$ for $\mathfrak{sl}_3$), but the composite-field OPE (involving the $\Lambda = {:}TT{:}$ composite) introduces double-pole Wick contractions that prevent naive reduction.

\item\label{nf:yangian-h4}
\emph{$\dim H^4(\barBgeom(Y(\mathfrak{sl}_2))) \stackrel{?}{=} 82$.}
Confirms the Yangian rationality (Conjecture~\ref{conj:yangian-bar-gf}).  The Yangian has $4$ weight-$1$ generators (the RTT matrix entries), giving chain groups of manageable size: $\dim\bar{B}^4 = 4^4 \cdot 6 = 1536$.  This is computationally accessible.
\end{enumerate}

\subsection{Medium-term challenges (new families)}

\begin{enumerate}[label=\textbf{MF-\arabic*.}, leftmargin=3em]
\item \emph{$\widehat{G}_{2,k}$ bar cohomology at degrees $2$ and~$3$.}
The first genuinely non-simply-laced computation.  The chain groups have size $14^2 \cdot 1 = 196$ at degree~$2$ and $14^3 \cdot 2 = 5488$ at degree~$3$.  The degree-$2$ value is feasible; degree-$3$ requires the full $G_2$ structure constants (which involve the non-trivially-normalized short roots).

\item \emph{$\widehat{\mathfrak{sl}}_4$ through degree~$3$.}
With $\dim\mathfrak{sl}_4 = 15$, the chain groups are $15, 225, 6750$.  Three data points would determine the unique rational GF of degree $\operatorname{rank} + 1 = 4$ (if the rank-plus-one pattern holds), and would test whether the DS-invariant discriminant is a cubic polynomial.

\item \emph{Combinatorial model for $\widehat{\mathfrak{sl}}_2$.}
The Virasoro algebra has an explicit Motzkin path model (Theorem~\ref{conj:motzkin-path-model}).  An analogous model for $\widehat{\mathfrak{sl}}_2$ --- giving a bijection between $H^n(\barBgeom(\widehat{\mathfrak{sl}}_2))$ and a family of $R(n+3)$ combinatorial objects --- would illuminate the representation-theoretic meaning of the Riordan numbers.  Candidates: non-crossing partitions without singletons on $\{1, \ldots, n+3\}$; restricted lattice paths with three step types and a non-touching constraint.
\end{enumerate}

\subsection{Long-range problems}

\begin{enumerate}[label=\textbf{LR-\arabic*.}, leftmargin=3em]
\item \emph{Intrinsic construction of the Kodaira--Spencer operator.}
Conjecture~\ref{conj:discriminant-ks-operator} predicts a rank-$(\operatorname{rank}\mathfrak{g}+1)$ transfer operator $\mathrm{KS}_{\mathfrak{g}}$ whose characteristic polynomial recovers the discriminant.  An intrinsic construction (from the Casimir operators on the bar complex, without first computing the Hilbert series) would reduce the bar cohomology computation to a finite-dimensional eigenvalue problem.

\item \emph{Uniformization of the algebraic hierarchy.}
Is there a structural reason why rank-$1$ algebras have algebraic (degree~$2$) bar GFs while rank-$2$ algebras have rational GFs?  A possible framework: the Picard--Fuchs D-module of the bar GF has holonomic rank $\operatorname{rank}\mathfrak{g} + 1$; for rank~$1$ this is a second-order Fuchsian ODE (with algebraic solutions), while for rank~$2$ it is third-order (and the solution space may contain only rational functions if the monodromy is trivial).

\item \emph{Non-simply-laced discriminant structure.}
For non-simply-laced $\mathfrak{g}$, the DS reduction $\widehat{\mathfrak{g}}_k \twoheadrightarrow \mathcal{W}^k(\mathfrak{g})$ involves the dual Coxeter number $h^\vee \neq h$.  The discriminant should reflect this discrepancy.  For $B_2 = \mathfrak{sp}_4$: $h = 4$, $h^\vee = 3$, $\dim = 10$.  Does $\Delta_{\mathfrak{sp}_4}$ factor over $\mathbb{Q}(\sqrt{d})$ for some discriminant~$d$ involving both~$h$ and~$h^\vee$?

\item \emph{Periodicity and the Coxeter number.}
The Kac--Moody periodicity theorem gives period $2h$ for $\widehat{\mathfrak{sl}}_2$ (the $E_2$-page has period~$4$ in conformal weight).  Whether this periodicity manifests as a functional equation for the bar GF --- of the form $P(x) = \text{rational} + x^{2h} P(x)$, analogous to a modular relation --- is unknown.

\item \emph{$W_\infty$ and infinite-rank limits.}
The MC4 tower $\mathcal{W}_N$ as $N \to \infty$ should converge to $W_{1+\infty}$.  The bar cohomology dimensions $\dim(\mathcal{W}_N^!)_n$ form a two-parameter array $(N, n)$.  Is there a closed-form bivariate generating function $\sum_{N,n} \dim(\mathcal{W}_N^!)_n\, t^N x^n$?  For fixed~$n$, the $N$-dependence is polynomial (the number of generators grows linearly in~$N$), so the bivariate GF should be rational in~$t$.
\end{enumerate}


\section{Summary of OEIS cross-references}
\label{sec:oeis}

\begin{center}
\renewcommand{\arraystretch}{1.4}
\begin{tabular}{lllp{5.5cm}}
\toprule
\textbf{OEIS} & \textbf{Sequence name} & \textbf{First terms} & \textbf{Role in the theory} \\
\midrule
A000041 & Partition numbers & 1,1,2,3,5,7,11 & Heisenberg/fermion bar cohomology \\
A000108 & Catalan numbers & 1,1,2,5,14,42 & Planar trees (HTT), associahedra \\
A001006 & Motzkin numbers & 1,1,2,4,9,21,51 & Virasoro bar $= M(n{+}1) - M(n)$ \\
A002026 & Motzkin differences & 0,1,2,5,12,30,76 & Virasoro bar cohomology directly \\
A005043 & Riordan numbers & 1,0,1,1,3,6,15,36 & $\widehat{\mathfrak{sl}}_2$ Koszul dual dims \\
A025565 & Central Delannoy (shifted) & 1,2,4,10,26,70 & $\beta\gamma$ bar cohomology \\
A132045 & $2^n - n + 1$ & 2,3,6,13,28,59 & $bc$ ghost bar cohomology \\
\bottomrule
\end{tabular}
\end{center}

\medskip

\noindent
The sequence $\dim(\mathcal{W}_3^!)_n = 2, 5, 16, 52, \ldots$ does not appear in OEIS as of this writing.  If the conjectured recurrence $a(n) = 3a(n-1) + a(n-2) - 1$ is correct, the sequence continues $171, 564, 1862, 6149, \ldots$ and would be a new entry.  Similarly, the $\widehat{\mathfrak{sl}}_3$ sequence $8, 36, 204, \ldots$ (if confirmed through higher degrees) would be new.

\medskip

\emph{A naturality observation.}  The known bar cohomology sequences --- partition numbers, Riordan numbers, Motzkin numbers, central Delannoy numbers --- all belong to the class of lattice path counts (with varying step sets and boundary conditions).  This is not a coincidence: the PBW spectral sequence, which computes bar cohomology via iterated Lie algebra cohomology, can be reformulated as a lattice path enumeration problem, where each conformal weight level corresponds to a horizontal coordinate and each cohomological degree to a vertical coordinate.  The step set is determined by the OPE singularity structure: simple poles give rise to ``up'' and ``down'' steps, higher-order poles to ``flat'' steps, and the boundary condition (non-negativity of height) reflects the unitarity of the underlying vertex algebra.  A precise formulation of this lattice-path / bar-cohomology correspondence, extending the Motzkin path model of Theorem~\ref{conj:motzkin-path-model} to all Koszul chiral algebras, would constitute a significant structural advance.
